%!TEX root = ../../csuthesis_research.tex

\chapter{研究计划、可行性分析与总结}

\section{研究计划与进度安排}

本课题进度安排如表~\ref{tab:plan} 所示。

\begin{table}[htbp]
\centering
\caption{课题研究进度计划（调研报告阶段制定）。}
\label{tab:plan}
\begin{tabular}{p{2.8cm}p{9.2cm}}
\toprule
时间阶段 & 主要工作内容 \\
\midrule
第1周 & 文献精读与问题定义，确定技术路线、符号系统与实验协议。 \\
第2周 & 完成统一代码骨架与数据处理流程，实现基础闭式解模块。 \\
第3周 & 完成疾病分类实验初版，输出首轮对比结果与误差分析。 \\
第4周 & 完成时间序列实验与关键消融，验证多尺度与高维映射增益。 \\
第5周 & 完成语义分割实验初版，重点评估旧类保持与语义漂移抑制。 \\
第6周 & 完成多模态路由实验，统计路由准确率与后向迁移指标。 \\
第7周 & 汇总跨任务结果，统一图表与结论，补充敏感性与稳定性分析。 \\
第8周 & 完成调研报告定稿、查重自检与汇报材料准备。 \\
\bottomrule
\end{tabular}
\end{table}

在上述周计划基础上，进一步给出阶段性交付物：
\begin{itemize}
    \item \textbf{第1--2周交付}：文献综述草稿、统一符号表、可运行代码骨架；
    \item \textbf{第3--4周交付}：分类与时序主结果、关键消融结果；
    \item \textbf{第5--6周交付}：分割与多模态主结果、对比分析文本；
    \item \textbf{第7--8周交付}：定稿调研报告、汇报 PPT、实验复现实验包。
\end{itemize}

为使各周交付物可量化追踪，进一步约定每周需提交统一日志模板与关键指标快照：分类与时序任务以 Acc / Forg 为周度核心指标，分割任务以新类 / 旧类 mIoU 为核心，多模态路由任务以 Avg-Task 与 BWT 为核心；同时辅以增量阶段耗时、显存峰值、$\Psi_t$ 条件数三类工程效率指标，作为周度回归检测的旁路证据。

\section{工作量评估与资源配置}

\subsection{工作量评估}

按章节与任务复杂度估算，工作量主要分布为：
\begin{itemize}
    \item 理论推导与证明整理：约 15\%（含递归岭回归闭式解、Woodbury 恒等式、隐私保护性质三类定理的整理与排版）；
    \item 四类任务实验实现与调参：约 55\%（分类约 10\%、时序约 12\%、分割约 20\%、多模态路由约 13\%，分割与多模态路由占比相对较高）；
    \item 图表制作与结果分析：约 20\%（含主结果表、任务演化曲线、消融表、效率图四类输出模板）；
    \item 论文写作与修改：约 10\%（含调研报告与毕业论文章节互通时的术语、记号与图表对齐工作）。
\end{itemize}

\subsection{资源配置}

建议采用统一代码框架、统一日志模板和统一绘图脚本，避免后期跨章节重复劳动。对于计算资源，优先保障语义分割与多模态实验时段，分类与时序实验可并行推进。

\section{可行性分析与风险应对}

\subsection{可行性基础}

课题具备较好的实施条件：
\begin{itemize}
    \item 理论可行：递归岭回归闭式解具有明确数学推导基础；
    \item 数据可行：四类任务均有公开基准可复现；
    \item 工程可行：方案以线性分类头更新为核心，计算与实现复杂度可控。
\end{itemize}



\section{本章小结与全文总结}

总览全文：第一章界定持续学习问题，并将研究背景从经典图像分类延伸到智能体记忆与自进化等近年快速演进的方向，明确闭式解路线对这些隐私敏感、长程累积场景的契合性；第二章构建递归岭回归闭式解的统一理论体系，给出与联合训练等价的递推公式、复杂度分析以及隐私保护性质；第三章从模块化角度提出统一方案，并将其与持续学习的数据 / 算力 / 稳定性 / 扩展性四类约束逐一对应；第四章将统一方案细化为医疗图像分类、时间序列分类、二维与三维语义分割、多模态大语言模型专家路由四类应用框架；第五章给出统一实验协议、评价指标、对比方法、消融与风险设计；第六章给出可执行的进度安排与可行性论证。这份调研报告为后续执行奠定了坚实基础。
